\documentclass[10pt]{article}

\usepackage[utf8]{inputenc}

\usepackage[T1]{fontenc}

\usepackage{amsmath}

\usepackage{amsfonts}

\usepackage{amssymb}

\usepackage[version=4]{mhchem}

\usepackage{stmaryrd}

\usepackage{graphicx}

\usepackage[export]{adjustbox}

\graphicspath{ {./images/} }

\usepackage{bbold}



\begin{document}

Ненаблюдаемость кварков в свободном состоянии (конфайнмент), а также противоречие с принципом Паули, возникающее при построении ряда адронов из кварков (напр., $\Omega^{-} \sim s s s$ ) и заключающееся в невозможности антисимметризовать волновую функщию, долгое время препятствовали признанию кварков как реальных объектов. Разрешение этих трудностей было найдено на основе гипотезы о существовании дополнительного квантового числа, принимающего три значения и названного цветом. Это утроило число состояний кварков, но число адронов осталось тем же, поскольку явление конфайнмента было сведено к постулату о физич. наблюдаемости только синглетных относительно группы цвета $S U(3)_{\text {color }}$ состояний.



В рамках К. м. с цветными кварками естественным образом были объяснены многие качественные закономерности сильных взаимодействий - такие, напр., как аннигиляционная модель слабых и электромагнитных взаимодействий, правило Окубо-Цвейга-Изуки, партонная картина глубоко-неупругого рассеяния и инклюзивных реакций, правила кваркового счета для асимптотич. поведения упругих формфакторов и др. Был вычислен также ряд количественных параметров, среди к-рых магнитные моменты барионов, ширина распада $\pi^{\circ} \rightarrow \gamma$, сечение $e^{-} e^{+}-$аннигиляции в адроны, массовые сдвиги внутри унитарных мультиплетов и др. В рамках К. м. на основе использования нерелятивистского подхода, приближения кваркового мешка и других были рассчитаны спектр масс низколежаших адронов и множество ширин перехода между ними. В этом подходе необходима такая динамич. характеристика кварков, как масса. Вследствие ненаблюдаемости кварков в свободном состоянии понятие массы кварка вводится в рамках квантовой хромодинамики как «токовая» или лагранжева. Отражением существования сушественно различных масштабов масс адронов $\left(m_{p} \sim 1\right.$ ГэВ, $m_{\gamma} \sim 10$ ГэВ) является большое различие в величинах токовых масс кварков $m_{u}=m_{d}=0$, $m_{s}==150 \mathrm{MэB}, m_{c}=1,3$ ГэВ, $m_{b}=4,5$ ГэВ, для массы экспериментально еще не обнаруженного $t$-кварка существует ограничение $m_{t}>60$ ГэВ.



Экспериментальное подтверждение наличия внутри адронов структурных единиц, обладающих именно теми свойствами, к-рые теория приписывает кваркам, позволяет говорить о существовании кварков как реальных объектов, из к-рых адроны действительно состоят, хотя полное теоретич. понимание механизма образования адронов из кварков в настоящее время отсутствует.



Модель цветных кварков, дополненная требованием локальной цветовой симметрии, служит основой квантовой хромодинамики - современной количественной теории сильных взаимодействий.



Лum.: [1] Боголюбов Н. Н., Струминский Б. В., Тавхе л и д з е А. Н., К вопросу о составных моделях в теории элементарных $\begin{array}{ll}\text { частиц, Дубна, } 1965 . & \text { A. А. Пивоваров. }\end{array}$



КВАРКОВЫЙ МЕШОК - феноменологическая модель $a \partial-$ ронов, позволяющая успешно рассчитывать их низкоэнергетические характеристики: спектр масс адронов, магнитные моменты, вершины нуклон-нуклонного и пион-нуклонного взаимодействий и т. д. В простейшем варианте (см. [1]) К. м. - это система эффективно свободных безмассовых фермионных полей с квантовыми числами валентных кварков, локализованных в конечной области пространства (напр., сферы радиуса $R$ ) и удовлетворяющих определенным граничным условиям.



Размер К. м. $R$ - внешний параметр, через к-рый выражаются все размерные характеристики адронов. В более сложных моделях К. м. [2], [3] предполагается существование отрицательной плотности энергии вакуума внутри мешка, что делает К. м. энергетически стабильным, по крайней мере для простейших адронных состояний. Удовлетворительно описывающим наблюдаемую $S U(2) \otimes S U(2)$-киральную симметрию ядерных сил является киральный К. м., в к-ром вводятся дополнительно безмассовые (в киральном пределе) пионные поля (см. [3]). Граничные условия на границе кирального К. м. связывают кварковые и пионные поля так, чтобы аксиальный ток, соответствующий киральной симметрии, сохранялся.



Jum.: [1] Bogolioubov N. N., "Ann. Inst. H. Poincaré», 1968, t. 8, p. 163-90; [2] Chodos A., Thorn G., «Phys. Lett.», 1974, v. B 53, p. 359-61; [3] B rown G., Rho M., там же, 1979, v. B82,



\begin{center}

\includegraphics[max width=\textwidth]{2023_07_08_8bac2044815b1cb06ec9g-1}

\end{center}



KBATEPHИОН - элемент множества $\mathbb{H}$, представимый в виде



\[

q=\alpha_{0} 1+\alpha_{1} i+\alpha_{2} j+\alpha_{3} k

\]



где $\alpha_{0}, \alpha_{1}, \alpha_{2}, \alpha_{3}-$ действительные числа, а $(1, i, j, k)-$ образуюшие базиса в $\mathbb{H}$, удовлетворяюшие соотношениям:



\[

\left.\begin{array}{c}

1 i=i, 1 j=j, 1 k=k, 1^{2}=1, i^{2}=j^{2}=k^{2}=-1, \\

j i=-j i=k, k i=-i k=j, j k=-k j=i .

\end{array}\right\}

\]



Умножение К. $q$ на скаляр $\alpha$ и сложение К. определяются так же, как и для обычных векторов. Можно ввести произведение двух К. $q=\alpha_{i} e_{i}$ и $q^{\prime}=\beta_{i} e_{i}$ формулой



\[

q q^{\prime}=\sum_{i, j} \alpha_{i} \beta_{j} e_{i} e_{j},

\]



иногда выделяют скалярную и векторную части К.:



\[

q=\alpha_{0}+V

\]



тогда умножение векторных частей определяется формулой



\[

V_{1} V_{2}=-\left(V_{1} V_{2}\right)+\left[V_{1} V_{2}\right]

\]



Тем самым множество $\mathbb{H}$ преврашается в алгебру (а л г е 6 р у к в а т е р н и о нов). Из соотношений (1) следует, что $\mathbb{H}-$ некоммутативная, но ассоциативная алгебра. Алгебра $\mathbb{H}$ содержит в виде подалгебры поле действительных чисел $\mathbb{R}=\left\{\alpha e_{0}\right\}$ и поле комплексных чисел $\mathbb{C}=\left\{\alpha e_{0}+\beta e_{1}\right\}$.



Алгебра $\mathbb{H}$ допускает изоморфное матричное представление с помошью Паули матриц:



\[

e_{0}=\left\|\begin{array}{ll}

1 & 0 \\

0 & 1

\end{array}\right\|, \quad e_{1}=i\|\| \begin{array}{ll}

0 & 1 \\

1 & 0

\end{array}\left\|, \quad e_{2}=i\right\| \begin{array}{cc}

0 & i \\

-i & 0

\end{array}\left\|, \quad e_{3}=i\right\| \begin{array}{cc}

1 & 0 \\

0 & -1

\end{array} \| .

\]



Для каждого К. $q=\alpha_{0} e_{0}+\alpha_{1} e_{1}+\alpha_{2} e_{2}+\alpha_{3} e_{3}$ определены сопряженный К. $\bar{q}=\alpha_{0} e_{0}-\alpha_{1} e_{1}-\alpha_{2} e_{2}-\alpha_{3} e_{3}$ и норма $N(q)=q \bar{q}=\bar{q} q=\alpha_{0}^{2}+\alpha_{1}^{2}+\alpha_{2}^{2}+\alpha_{3}^{2}=|q|^{2}$. Обратным кватернионом является $q^{-1}=\bar{q} /|q|^{2}$. Каждый ненулевой К. имеет обратный.



Наряду с действительными и комплексными числами в различных вопросах теории представлений групп, топологии и физики можно использовать К. Врашение трехмерного пространства можно задать при помощи К. с нормой 1 (аналогично тому, как вращение плоскости задается комплексным числом с модулем 1).



КЕЛЬВИНА ПРЕОБРАЗОВАНИЕ - преобразование функций, определенных в областях евклидова пространства $\mathbb{R}^{n}$, $n \geqslant 3$, при котором гармонические функции переходят в гармонические. Получено У. Томсоном (W. Thomson, 1847) (лордом Кельвином).



Если $u(x)$ - гармонич. функция в области $D \subset \mathbb{R}^{n}$, то ее К. П. есть функция



\[

v(y)=\left(\frac{R}{|y|}\right)^{n-2} u\left(\frac{R^{2}}{|y|^{2}}, y\right), \quad v(\infty)=0,

\]



гармоническая в области $D^{*}$, получающейся из $D$ инверсией относительно сферы $S_{R}=\{x:|x|=R\}$, то есть отображением пространства $\mathbb{R}^{n}$, определяемым формулами



\[

x \rightarrow y=\frac{R^{2}}{|x|^{2}} x, \quad 0 \rightarrow \infty

\]



где $x=\left(x_{1}, \ldots, x_{n}\right), \quad|x|=\left(x_{1}^{2}+\ldots+x_{n}^{2}\right)^{1 / 2}$.



При К. П. гармонич. функции $u(x)$ в областях $D$, содержаших $\infty$, регулярные в $\infty$, то есть такие, что



\[

\lim _{|x| \rightarrow \infty} u(x)=0

\]





\end{document}